\section{Completion and the lattice consequence}
\label{sec:completion}

\begin{proof}[Proof of Theorem~\ref{thm:main}]
The reduction in Section~\ref{sec:preliminaries} permits us to work
with a simple graph. By Lemma~\ref{lem:orbit-growth}, its uniform
exponential growth rate exists. If this rate is greater than one,
Theorem~\ref{thm:exponential} applies. Otherwise the graph has
subexponential growth.

In this remaining case, if \eqref{eq:superpoly-growth} holds,
Theorem~\ref{thm:superpoly} excludes critical percolation. If it does
not hold, \eqref{eq:polynomial-subsequence} holds along an unbounded
sequence. Proposition~\ref{prop:nilpotent-model} supplies the group
action required by Theorem~\ref{thm:structured}, which again excludes
critical percolation when $\pc<1$. Thus every fixed vertex has
probability zero of belonging to an infinite critical cluster.
The countable union over vertices proves the almost-sure assertion.
\end{proof}

\begin{corollary}
\label{cor:lattices}
For nearest-neighbor Bernoulli bond percolation on $\Z^d$, for every
$d\ge2$, there is almost surely no infinite cluster at the critical
probability. In particular this holds on $\Z^3$.
\end{corollary}

\begin{proof}
The graph is infinite, connected, locally finite, and transitive.
We verify $\pc<1$ by a planar boundary count on a coordinate copy
of $\Z^2$. If its origin cluster is finite, the dual edges crossing
its edge boundary contain a simple cycle surrounding the origin;
all crossed primal edges are closed. To see the existence of such
a cycle, the dual boundary has even degree at every vertex and
decomposes into simple cycles. The positive horizontal ray from
the origin crosses the whole boundary an odd number of times, so
one cycle crosses this ray oddly. A closed cycle crosses the whole
horizontal line evenly, and hence also crosses the negative ray.

A cycle of length $n$ therefore has a positive-ray crossing at
distance less than $n$ from the origin. Choosing that crossing,
an orientation, and a nonbacktracking continuation bounds the
number of candidates by $2n3^n$. Independence bounds the probability
that the origin cluster is finite by
\[
 \sum_{n\ge1}2n3^n(1-p)^n.
\]
For $p<1$ sufficiently close to one this sum is less than one,
so planar percolation occurs with positive probability. The same
configuration is a subgraph of $\Z^d$, proving $\pc(\Z^d)<1$.
Theorem~\ref{thm:main} now applies.
\end{proof}
